%==============================================================================
% Research Methods - Group assignment - proposals
%==============================================================================

% This loads the article template. For a text written in English, add the
% "english" option, i.e. \documentclass[english]{hogent-article}
\documentclass{hogent-article}

% Include bibliography file
\addbibresource{references.bib}

% Info about the assignment, course and study programme

\studyprogramme{Professionele bachelor toegepaste informatica}
% TODO: replace 2N1 with the correct class group(s)!
\course{Research Methods, 2N1} 
\assignmenttype{Groepsopdracht}
\academicyear{2025-2026}
\title{Schriftelijke voorbereiding onderzoeksvoorstellen}

% TODO: English text? Comment or remove the lines above and uncomment the lines
% below

% \studyprogramme{Professional bachelor in applied informatics}
% \course{Research Methods, INT-IT/2B}
% \assignmenttype{Group assignment}
% \academicyear{2025-2026}
% \title{Written preparation of research proposals}

\author{Simon Borghart}
\email{simon.borghart@student.hogent.be}

\author{Emiel De Pelsmaeker}
\email{emiel.depelsmaeker@student.hogent.be}

\author{Simon Boey}
\email{simon.boey@student.hogent.be}

\projectrepo{https://github.com/SaZoMi/RM-Group-64-Specialisatie.git}

\specialisation{AI \& Data Engineering}

\keywords{large language models, AI-agenten, kostenoptimalisatie, caching, prompt-compressie, model-routing}

\begin{document}

\begin{abstract}
  Bestaande technieken voor het cachen van LLM-antwoorden, zoals context caching en semantische caching, zijn ontworpen voor chatbot-achtige toepassingen en schieten tekort bij AI-agenttoepassingen. Agenttoepassingen voeren immers meerdere planningsstappen na elkaar uit en hun output hangt af van externe data en omgevingscontext, waardoor klassieke caching daar beperkt bruikbaar is \autocite{Zhang2025}. Dit zorgt voor onnodig hoge operationele kosten en latency bij bedrijven die LLM-agenten inzetten in productie. Dit onderzoeksvoorstel beschrijft een geplande studie naar de vraag welke combinatie van kostenoptimalisatietechnieken het meest geschikt is voor LLM-agenttoepassingen, waarbij naast caching ook prompt-compressie en model-routing onderzocht worden. Op basis van een eerste literatuurstudie wordt een probleem- en oplossingsdomein afgebakend. Daarna volgt een voorstel van methodologie waarin de technieken vergelijkend geëvalueerd worden op kosten- en kwaliteitscriteria, specifiek voor agentgebaseerde toepassingen. Het onderzoek moet resulteren in een onderbouwd handvat voor IT-professionals om LLM-kosten te beheersen bij het ontwerpen van AI-agenttoepassingen.
\end{abstract}

\tableofcontents

\section{Inleiding}%
\label{sec:inleiding}

Het aantal toepassingen dat gebruikmaakt van large language models (LLM's) via commerciële API's, zoals die van OpenAI en Anthropic, of via zelf gehoste modellen, is de afgelopen jaren sterk toegenomen. Een groeiend deel van deze toepassingen zijn zogenaamde AI-agenten: systemen die zelfstandig een taak plannen, tussenstappen uitvoeren en hierbij herhaaldelijk een LLM aanroepen. Dit onderzoeksvoorstel richt zich specifiek op de operationele kostenproblematiek van dergelijke agenttoepassingen, en niet op de kwaliteit van de gegenereerde output op zich of op ethische aspecten van AI-gebruik.

Zhang et al. (2025) tonen aan dat bestaande LLM-cachingtechnieken, zoals context caching en semantische caching, primair ontworpen zijn voor het bedienen van chatbots \autocite{Zhang2025}. Bij chatbots is de output doorgaans een direct antwoord op een vraag, waardoor gelijkaardige of identieke vragen relatief eenvoudig herkend en uit een cache beantwoord kunnen worden. Bij AI-agenten ligt dit anders: de output van een planningsstap hangt af van externe data en van de specifieke omgevingscontext op dat moment, waardoor twee semantisch gelijkaardige taken toch tot verschillende, niet-herbruikbare plannen kunnen leiden. Het gevolg is dat klassieke caching bij agenttoepassingen minder herbruikbaarheid oplevert dan bij chatbots, en dat de kosten en latency van agenttoepassingen daardoor onnodig hoog blijven \autocite{Zhang2025}.

Daarnaast wordt in de vakliteratuur op technieken als prompt-compressie en model-routing gewezen als aanvullende manieren om kosten te verlagen \autocite{Chen2024}. Het ontbreekt echter aan een overzicht dat deze technieken specifiek afweegt in de context van agenttoepassingen, waar de klassieke aannames van caching niet zonder meer opgaan. Daardoor hebben IT-professionals die AI-agenten ontwikkelen momenteel weinig onderbouwde houvast bij het maken van een keuze.

Dit leidt tot de volgende hoofdonderzoeksvraag: \emph{Welke combinatie van kostenoptimalisatietechnieken (agentgeschikte caching, prompt-compressie en model-routing) is het meest geschikt om de operationele kosten van LLM-agenttoepassingen in een productieomgeving te verlagen, zonder de kwaliteit van de output significant te verminderen?}

Deze hoofdvraag wordt opgesplitst in deelvragen binnen het probleemdomein en het oplossingsdomein.

Deelvragen probleemdomein:
\begin{itemize}
  \item Waarom schieten klassieke context- en semantische caching specifiek tekort bij LLM-agenttoepassingen?
  \item Welke factoren bepalen de kostenstructuur van LLM-agenttoepassingen, bijvoorbeeld het aantal planningsstappen, promptlengte en modelkeuze per stap?
  \item Welke risico's ontstaan wanneer kosten worden verlaagd zonder rekening te houden met de specifieke werking van agenttoepassingen, bijvoorbeeld op vlak van kwaliteitsverlies of onbetrouwbare planning?
\end{itemize}

Deelvragen oplossingsdomein:
\begin{itemize}
  \item Hoe werkt agentgeschikte caching, zoals het cachen en hergebruiken van planstructuren in plaats van volledige antwoorden, en in welke situaties levert dit meetbare kostenbesparing op?
  \item Wat houdt prompt-compressie in en in welke mate is deze techniek toepasbaar op de planningsstappen van een agent?
  \item Hoe werkt model-routing, waarbij dynamisch een goedkoper of kleiner model wordt gekozen voor eenvoudigere planningsstappen, en welke tools of frameworks ondersteunen dit?
  \item Hoe kan de impact van deze technieken op kosten en kwaliteit gemeten en onderling vergeleken worden bij agenttoepassingen?
\end{itemize}

Het doel van dit onderzoek is een onderbouwd en vergelijkend overzicht te bieden van kostenoptimalisatietechnieken voor LLM-agenttoepassingen in productieomgevingen, gericht op IT-professionals die dergelijke toepassingen ontwerpen of beheren. De rest van dit voorstel is als volgt opgebouwd: sectie~\ref{sec:literatuurstudie} bespreekt de belangrijkste inzichten uit de literatuur, sectie~\ref{sec:methodologie} beschrijft de voorgestelde onderzoeksmethodologie, sectie~\ref{sec:resultaten} schetst de verwachte resultaten, en sectie~\ref{sec:conclusie} vat het voorstel samen.

\section{Literatuurstudie}%
\label{sec:literatuurstudie}

Zhang et al. (2025) benoemen het kernprobleem waarop dit voorstel is gebaseerd: klassieke context caching en semantische caching zijn ontwikkeld voor chatbot-toepassingen en houden geen rekening met het feit dat de output van een AI-agent afhangt van externe data en omgevingscontext \autocite{Zhang2025}. Als oplossing stellen zij Agentic Plan Caching (APC) voor, waarbij niet het volledige antwoord, maar een herbruikbare planstructuur wordt gecachet en aangepast aan de specifieke taak. Hun evaluatie op meerdere praktijktoepassingen toont een gemiddelde kostenreductie van 50,31\% en een latencyreductie van 27,28\%, bij behoud van de oorspronkelijke prestaties \autocite{Zhang2025}. Dit resultaat bevestigt zowel het bestaan van het probleem als de haalbaarheid van een oplossing die specifiek rekening houdt met de werking van agenten.

De bredere literatuur over kostenreductie bij LLM-gebruik onderscheidt daarnaast strategieën als prompt-adaptatie, modelbenadering en model-cascades \autocite{Chen2024}. Chen et al. (2024) tonen met FrugalGPT aan dat een combinatie van deze strategieën tot 98\% kostenreductie kan opleveren bij een vergelijkbare kwaliteit als het duurste beschikbare model. Deze bevindingen zijn echter voornamelijk getoetst op enkelvoudige LLM-aanroepen en niet specifiek op de meerstaps-planningsstructuur van agenttoepassingen.

Op het vlak van caching onderzoeken Regmi en Pun (2024) semantische caching. Hierbij worden niet enkel identieke, maar ook vergelijkbare vragen op basis van semantische gelijkenis uit een cache beantwoord \autocite{Regmi2024}. Liu et al. (2025) reiken hiervoor een theoretisch onderbouwd raamwerk aan voor cache-eviction onder onzekere query- en kostenverdelingen \autocite{Liu2025}. Beide werken bevestigen, net als Zhang et al. (2025), dat semantische caching in de basisvorm vooral geschikt is voor vraag-antwoordpatronen zoals bij chatbots.

Voor prompt-compressie blijft LLMLingua van Jiang et al. (2023) een veelgeciteerde referentie. Door minder relevante tokens uit de prompt te verwijderen, wordt de invoer verkleind met behoud van de essentiële betekenis \autocite{Jiang2023}. Deze techniek is in principe ook toepasbaar op de tussentijdse planningsprompts van een agent.

Op het vlak van model-routing bieden Dekoninck et al. (2024) een verenigd theoretisch kader dat routing, waarbij één model per vraag gekozen wordt, en cascading, waarbij opeenvolgend grotere modellen worden aangesproken tot een bevredigend antwoord, combineert tot een geoptimaliseerde strategie \autocite{Dekoninck2024}. Zhen et al. (2025) plaatsen deze technieken in een breder overzicht van efficiënte LLM-inferentie, met aandacht voor zowel technieken op instantieniveau als op clusterniveau \autocite{Zhen2025}.

Uit deze literatuur blijkt dat het probleem van beperkte cachebaarheid bij agenttoepassingen recent en concreet aangetoond is, en dat prompt-compressie en model-routing potentieel aanvullende oplossingen bieden. Vergelijkend onderzoek dat deze drie technieken samen afweegt specifiek voor agenttoepassingen is echter nog beperkt beschikbaar.

\section{Methodologie}%
\label{sec:methodologie}

Om de onderzoeksvragen te beantwoorden wordt een gefaseerde aanpak voorgesteld.

In een eerste fase, gericht op het probleemdomein, wordt via literatuuronderzoek geanalyseerd waarom klassieke caching tekortschiet bij agenttoepassingen en welke factoren de kostenstructuur van dergelijke toepassingen concreet bepalen.

In een tweede fase, gericht op het oplossingsdomein, worden agentgeschikte caching, prompt-compressie en model-routing elk afzonderlijk beschreven aan de hand van de geraadpleegde vakliteratuur. Deze beschrijving wordt aangevuld met een vergelijkende analyse op basis van gepubliceerde experimentele resultaten over kostenbesparing, kwaliteitsbehoud en impact op latency, telkens getoetst aan de specifieke context van agenttoepassingen.

In een derde fase worden de bevindingen samengebracht in een afwegingskader dat aangeeft in welke situaties welke techniek of combinatie van technieken aangewezen is voor LLM-agenttoepassingen.

\section{Verwachte resultaten}%
\label{sec:resultaten}

Het onderzoek moet een vergelijkend overzicht opleveren van kostenoptimalisatietechnieken voor LLM-agenttoepassingen. Voor elke techniek wordt een inschatting verwacht van de haalbare kostenbesparing, de invloed op kwaliteit en latency, en de situaties waarin de techniek het meest geschikt is. Op basis hiervan wordt een afwegingskader verwacht dat IT-professionals kan ondersteunen bij het maken van een onderbouwde keuze bij het ontwerpen van agenttoepassingen.

\section{Discussie en conclusie}%
\label{sec:conclusie}

Dit voorstel kadert een concreet en recent aangetoond probleem: klassieke LLM-cachingtechnieken schieten tekort bij AI-agenttoepassingen, waardoor de operationele kosten en latency van dergelijke toepassingen onnodig hoog blijven. Op basis hiervan wordt een onderzoek voorgesteld naar de effectiviteit van agentgeschikte caching, prompt-compressie en model-routing als oplossingsstrategieën. De voorgestelde methodologie combineert literatuuronderzoek met een vergelijkende analyse, met als doel een praktijkgericht afwegingskader op te leveren. Een mogelijke beperking van het onderzoek is de snelheid waarmee dit vakgebied evolueert, waardoor specifieke kostencijfers snel kunnen verouderen. Het voorgestelde afwegingskader is daarom opgezet om ook op langere termijn bruikbaar te blijven.

\printbibliography[heading=bibintoc]

\end{document}